\subsection{One contour argument for the two lattices}
\label{ca:lattice-method}

The internal wire sum and the later color Fourier sum have the same
analytic form. We give the contour argument once. It explains why a
Gaussian critical point contributes to these finite sums with the
claimed coefficient.

\begin{latticelemma}
Let $\mathcal B\subset\mathbb R^k$ be a fixed rectangle and let
$F(y,t)$ and $b(y,t)$ be holomorphic on a fixed complex neighborhood
of $\mathcal B\times\{0\}$, where $y\in\mathbb C^k$ and
$t\in\mathbb C^s$. Fix a real number $\rho$. Suppose that $F(y,0)$ is real for real
$y$, has a unique maximum $y_0$ in the interior of $\mathcal B$,
and has negative definite Hessian there. Suppose also that
$b(y_0,0)\ne0$.

Let $b_N=b+O(N^{-1})$ holomorphically on this neighborhood, with
the same bound for every fixed number of derivatives on smaller
neighborhoods. Then, for small complex $t$,
\begin{align*}
 &\sum_{j/N\in\mathcal B}
       N^\rho b_N(j/N,t)e^{NF(j/N,t)}\\
 &\qquad=
 N^{\rho+k/2}
 \frac{(2\pi)^{k/2}b(y_*(t),t)}
      {\sqrt{\det[-F_{yy}(y_*(t),t)]}}
 e^{NF(y_*(t),t)}(1+O(N^{-1})).
\end{align*}
The sum is over $j\in\mathbb Z^k$. Here $y_*(t)$ is the
continuation of $y_0$, and the square root is positive at $t=0$.
Precisely, after division by the displayed power, amplitude and
exponential, the remainder is a holomorphic function $R_N(t)$ with
$R_N=O(N^{-1})$. Each fixed derivative of $R_N$ has the same bound
on smaller parameter polydisks. The lattice is $N^{-1}\mathbb Z^k$;
its fundamental cell has volume $N^{-k}$.
\end{latticelemma}

\begin{proof}
The implicit-function theorem gives $y_*(t)$ and its nonsingular
Hessian. At $t=0$, compactness gives a fixed phase loss on the
boundary of $\mathcal B$ and outside a small ball about $y_0$.

In one coordinate, move each endpoint by at most $1/N$ to a
half-grid point without changing the enclosed lattice points.
The function $\pi N\cot(\pi Nz)$ has residue one at every
$j/N$. Integrate it around the resulting rectangle with horizontal
sides at heights $\pm a$, where $a>0$ is fixed and small.
On a vertical half-grid side its modulus is at most $\pi N$;
on the horizontal sides it is
$\mp\pi N\ii+O(Ne^{-2\pi Na})$.

Because the complete phase is real on the real rectangle, moving
an imaginary distance at most $a$ changes its real part by at most
$Ca^2$. The color perturbation costs at most $C\|t\|$ on this
fixed tube. Choose $a$ first, so that $Ca^2$ is smaller than one
quarter of the boundary loss and smaller than $\pi a$. Then shrink
the parameter polydisk. Every vertical side and every nonconstant
cotangent term has a fixed exponential loss relative to the
continued critical value.

The two constant horizontal terms each give one half of $N$ times
the real integral. Their connecting vertical segments have the
same boundary loss. Apply this argument in every coordinate.
The finite sum is therefore $N^k$ times the integral over
$\mathcal B$, up to an exponentially smaller absolute error.
The $O(N^{-1})$ strips between the half-grid rectangle and
$\mathcal B$ have the same fixed boundary loss. In particular there
is no multiplicity $2^k$.

Translate the real rectangle by $\ii\Im y_*(t)$. The translated
rectangle contains the critical point because $\Re y_*(t)$ stays
in its interior. Cauchy's theorem expresses the difference of the two
integrals as integrals over connecting faces, which retain the boundary
loss. On a small fixed ball the real part of $-F_{yy}$ remains positive
definite.
With $u=\sqrt N(y-y_*)$, the quadratic term gives the displayed
Gaussian determinant. The cubic phase term and linear amplitude
term are odd and integrate to zero. Taylor's formula and Gaussian
domination bound the remaining relative error by $O(N^{-1})$;
outside the ball there is a strict exponential loss. The factors
$N^k$ from the grid and $N^{-k/2}$ from the Gaussian give the
stated power. Holomorphy and Cauchy's formula give the derivative
bounds.
\end{proof}

The important distinction is between an exact finite sum and its local
Gaussian approximation. The cotangent integral first relates the
entire finite sum to a contour through the critical point, with an
exponentially smaller error. Only then is it legitimate to use the
Gaussian coefficient. The same argument will be applied to the
internal polynomial indices and to the auxiliary colors.
